\section{Soundness}
\label{sec:fol-deduction-03}

\begin{theorem}[Soundness of first-order natural deduction]
\label{thm:fol-soundness}
For every set \(\Gamma\cup\set{\varphi}\subseteq\Fm(\Lang)\),
\[
 \Gamma\derives_{\ND_{\FOL}}\varphi
 \quad\Longrightarrow\quad
 \Gamma\entails\varphi.
\]
\end{theorem}
\begin{proof}
Induct on the construction of the derivation, using the strengthened statement
that any structure and assignment satisfying all open assumptions satisfies
the conclusion.  The propositional cases are those of
\cref{thm:prop-nd-soundness}.

For \(\forall E\), if \(\Meta,s\entails\forall x\,\varphi\), then it
satisfies \(\varphi\) under
\(s[x\mapsto\sem{t}^{\Meta}_s]\).  The substitution lemma gives
\(\Meta,s\entails\varphi[t/x]\).  The \(\exists I\) case reverses this
argument: a satisfying value \(\sem{t}^{\Meta}_s\) supplies the existential
witness.

For \(\forall I\), let \(\Meta,s\) satisfy all open assumptions and choose an
arbitrary \(a\in\dom\Meta\).  Because \(x\) is free in no open assumption,
formula coincidence shows that \(s[x\mapsto a]\) also satisfies every open
assumption.  The induction hypothesis gives
\(\Meta,s[x\mapsto a]\entails\varphi\).  Since \(a\) was arbitrary,
\(\Meta,s\entails\forall x\,\varphi\).

For \(\exists E\), let \(\Meta,s\) satisfy the surviving open assumptions and
\(\exists x\,\varphi\).  Choose \(a\) such that
\(\Meta,s[x\mapsto a]\entails\varphi\), and put
\(s'=s[y\mapsto a]\).  Freshness of \(y\), formula coincidence, and the
substitution lemma give \(\Meta,s'\entails\varphi[y/x]\), while freshness also
ensures that \(s'\) satisfies every other open assumption.  The induction
hypothesis for the subsidiary derivation yields \(\Meta,s'\entails\psi\).
Because \(y\notin\FV(\psi)\), formula coincidence gives
\(\Meta,s\entails\psi\).

The \(\alpha\) rule is sound by alpha-invariance.  For \(=I\), actual
identity is reflexive.  For \(=E\), satisfaction of
\(r=t\) gives
\(\sem{r}^{\Meta}_s=\sem{t}^{\Meta}_s\).  The substitution lemma says that
the two instances of \(\varphi\) are evaluated under updates of \(s\) by
these equal elements, so satisfaction of one instance implies satisfaction of
the other.  This treats every rule of \(\ND_{\FOL}\) that is not a rule
of \(\ND_{\CPL}\), completing the induction.
\end{proof}

\begin{theorem}[Soundness of the first-order sequent calculus]
\label{thm:fol-lk-soundness}
Every sequent derivable in \(\LK_{\FOL}\) is satisfied by every
\(\Lang\)-structure and assignment.
\end{theorem}
\begin{proof}
Extend the induction in \cref{thm:prop-lk-soundness}: for each rule,
if every premise is satisfied by every structure and assignment, so is the
conclusion.  The rules \(\forall L\) and \(\exists R\) use the term value
and the substitution lemma.

For \(\forall R\), suppose toward a contradiction that some \(\Meta,s\)
does not satisfy the conclusion
\(\Gamma\Rightarrow\Delta,\forall x\,\varphi\) while the premise
\(\Gamma\Rightarrow\Delta,\varphi[y/x]\) is satisfied by every structure
and assignment.  Then \(\Meta,s\) satisfies every formula in
\(\operatorname{supp}(\Gamma)\), no formula in
\(\operatorname{supp}(\Delta)\), and not \(\forall x\,\varphi\).  The
universal satisfaction clause supplies an element \(a\in\dom\Meta\)
with \(\Meta,s[x\mapsto a]\notentails\varphi\).  Put
\(s'=s[y\mapsto a]\), where \(y\) is the eigenvariable.

First, \(\Meta,s'\notentails\varphi[y/x]\).  The substitution lemma
(\cref{thm:fol-substitution}) gives
\(\Meta,s'\entails\varphi[y/x]\) iff
\(\Meta,s'[x\mapsto\sem{y}^{\Meta}_{s'}]\entails\varphi\), and
\(\sem{y}^{\Meta}_{s'}=s'(y)=a\).  If \(y=x\), then
\(s'[x\mapsto a]=s[x\mapsto a]\) and the right side fails directly.  If
\(y\ne x\), then the eigenvariable condition---\(y\) occurs free in no
formula of \(\Gamma,\Delta,\forall x\,\varphi\)---gives
\(y\notin\FV(\varphi)\), so \(s'[x\mapsto a]\) and
\(s[x\mapsto a]\) agree on \(\FV(\varphi)\), and formula coincidence
(\cref{lem:fol-coincidence}) reduces the right side to
\(\Meta,s[x\mapsto a]\entails\varphi\), which fails.

Second, the same eigenvariable condition puts \(y\) outside the free
variables of every formula of \(\Gamma\) and \(\Delta\), so \(s\) and
\(s'\) agree on those free variables, and formula coincidence transfers
the pattern above to \(s'\): \(\Meta,s'\) satisfies every formula in
\(\operatorname{supp}(\Gamma)\) and no formula in
\(\operatorname{supp}(\Delta)\).  Together with
\(\Meta,s'\notentails\varphi[y/x]\), this means \(\Meta,s'\) does not
satisfy the premise, a contradiction.  The eigenvariable condition is what
licenses both uses of formula coincidence; without it, updating \(y\)
could change the truth values of the contexts or of \(\varphi\) itself.

The \(\exists L\) case is dual.  The equality initial sequent
\(\Rightarrow t=t\) is satisfied by every structure and assignment
because actual identity is reflexive, and \(=L\) preserves satisfaction by
every structure and assignment by the same equal-update argument used for
\(=E\) in \cref{thm:fol-soundness}.
\end{proof}

